\printedpage{119}
% head: A. Grothendieck : Dérivateurs, Ch. VI, Édition des Universités de Montpellier II et Paris VII
% foot: Version 9 novembre 2001

% the following paragraph is marked by a triple vertical rule in the left margin
\noindent\llap{$|||$\quad}Let $X \in \mathrm{Cat}/S$, $S = \begin{tikzcd}[row sep=tiny, column sep=small] & a \\ 0 \arrow[ur] \arrow[dr] & \\ & b \end{tikzcd}$. If $\Phi(X)$ is proper over $S$, then $X = \emptyset$ (if we assume $W \neq \mathrm{Fl}(\mathrm{Cat})$).

But along the way a bizarre idea appeared. Consider the infinite diagram, for $X$ in $\mathrm{Cat}/S$:
\[
X \longrightarrow \Phi(X) \longrightarrow \Psi\Phi(X) \longrightarrow \Phi\Psi\Phi(X) \longrightarrow \cdots
\]
If $X$ is arbitrary, these arrows are alternately in $W_S^{\mathrm{d}}$ and $W_S^{\mathrm{g}}$, so their composites are in $W_S$ and no more. But if $X$ is in $\und{Fib}_W S$, since this category is \textit{stable} under $\Phi$, $\Psi$, and since on $\und{Fib}_W S$ $W_S^{\mathrm{d}} = W_S^{\mathrm{g}} = W_S^{\mathrm{f}} = W_S^{\mathrm{univ}}$, we see that in this case the arrows are all in $W_S^{\mathrm{f}}$. We now feel like setting
\[
\Phi_\infty(X) = \varinjlim \text{ of all that}
\]
\msp{page 193}
and we would define symmetrically
\[
\Psi_\infty(X) = \varinjlim (X \longrightarrow \Psi(X) \longrightarrow \Phi\Psi(X) \longrightarrow \cdots).
\]

For this to be usable, we must postulate the axiom \underline{W(6 bis)} (compare XII, pages 14, 15). The direct limit of a sequence of morphisms $f_i : X_i \longrightarrow Y_i$, for a morphism of direct systems in Cat indexed by $\mathbf{N}$, is in $W$ if the $f_i$ are.

This axiom is close to W(6) of loc. cit., but is implied by it only with W(4b). It is stronger in the case of direct systems indexed by $\mathbf{N}$ (so also, more generally, of direct systems on filtering ordered sets having a countable cofinal set), but says nothing about direct systems on more general ordered index sets.

This axiom also implies that a filtering direct limit of morphisms $f_i$ is $W$-smooth, resp. $W$-proper, if the $f_i$ are. We then find that $\Phi_\infty(X)$ is both $W$-smooth (as a direct limit of $\Phi(X)$, $\Phi\Psi\Phi(X)$, \ldots, $\Phi(\Psi\Phi)^n(X)$, which are $W$-smooth over $S$) and $W$-proper (as a direct limit of $\Psi\Phi(X)$, \ldots, $(\Psi\Phi)^n(X)$ \ldots). On the other hand, that $X \longrightarrow \Phi_\infty(X)$ is in $W_S^{\mathrm{f}}$ - for $W_S^{\mathrm{f}}$ is stable under countable filtering direct limits in $\mathrm{Cat}/S$. We have thus found a functor
\msp{page 194}
\[
\Phi_\infty : \und{Fib}_W S \longrightarrow \und{Parf}_W S \subseteq \und{Fib}_W S
\]
and a functorial morphism
\[
i_X : X \longrightarrow \Phi_\infty(X), \qquad \boxed{i_X \in W_S^{\mathrm{f}}} \qquad \text{if } X \in \und{Fib}_W S.
\]
% the box around i_X \in W_S^f is printed with rounded corners

This functor then allows us, in the absence of W(4b), to do what the functor $C$ absolutely refused. It is true that $C$ has the advantage of being defined on all of $\mathrm{Cat}/S$, and of being
\printedpage{120}
% head: A. Grothendieck : Dérivateurs, Ch. VI, Édition des Universités de Montpellier II et Paris VII
% foot: Version 9 novembre 2001
reasonable and even very useful there, when we have W(4b). But $\Phi_\infty$ too is defined on all of $\mathrm{Cat}/S$, as is $i_X$
\[
\begin{aligned}
&\Phi_\infty : \mathrm{Cat}/S \longrightarrow \und{Parf}_W S\\
&i : \mathrm{id} \longrightarrow \Phi_\infty \qquad \text{i.e.}\ \begin{cases} i_X : X \longrightarrow \Phi_\infty(X)\\ i_X \in W_S, \end{cases}
\end{aligned}
\]
with the sole handicap (inevitable) that if $X$ is not in $\und{Fib}_W S$, we can say no better than $i_X \in W_S$, to the exclusion of $i_X \in W_S^{\mathrm{f}}$. But for $C(X)$ it was necessarily the same. So we get the impression that $\Phi_\infty$, take your pick, $\Psi_\infty$, six of one, half a dozen of the other, plays in the presence of W(6 bis) the same role as C in the presence of W(4b). At most (in the presence of both axioms at once) the functor $C$ has the aesthetic advantage of taking its values in $\und{Parf}_\omega S$, instead of merely $\und{Parf}_W S$ (\footnote[93]{\textit{In fact} $\Phi_\infty = C$, if $C$ is defined simply via one-sided infinite paths. If we define $C$ via \textit{finite} paths of all types $c$, then it does indeed take its values in $\und{Parf}_\omega S$.}). But it is imperative to % [sic] "que C en présence" (C in roman type) -> $C$, italic everywhere else
\msp{page 195}
realize now that $\Phi_\infty$ has an advantage of exactly the same order over $C$, by introducing
\[
W''_\omega = \text{smallest subset } \Sigma \text{ of } \mathrm{Fl}(\mathrm{Cat}) \text{ satisfying } \mathrm{W}(1, 2\text{ bis}, 3, 6\text{ bis})
\]
and note that
\[
\Phi_\infty(X) \in \mathrm{Ob}\,\und{Parf}_{W''_\omega} S, \qquad i_x \in W''_\omega % [sic] i_x -> i_X, the morphism is i_X everywhere else
\]
- but here I notice that we have
\[
W_\omega \subseteq W''_\omega,
\]
so $C$ is still a bit better in this respect \ldots

When we have $\Phi_\infty$ with the properties we have just made explicit, we find that the inclusion
\[
\und{Parf}_W S \longrightarrow \und{Fib}_W(S)
\]
induces an equivalence not only for the categories localized for $W_S$, but also for the fine localized categories
\[
(\und{Parf}_W S)(W_S^{\mathrm{f}})^{-1} \longrightarrow (\und{Fib}_W(S))(W_S^{\mathrm{f}})^{-1} \overset{\text{def}}{=} \mathrm{HOT}^{\mathrm{lc}}_W(S)
\]
(\footnote[94]{We can even replace $\und{Parf}_W S$ there by $\und{Parf}_{W''_\omega} S$, \textit{or} by any other full subcategory of $\und{Fib}_W S$ containing $\und{Parf}_{W''_\omega} S$.}). But we cannot for all that reduce the number of useful localized categories to three, in the absence of W(4b). At most, having a functor $\Phi_\infty$ which is simpler than $C$, in the absence of W(4b), because we know in any case that
\[
i_X : X \longrightarrow \Phi_\infty(X)
\]
is in $W_S^{\mathrm{f}}$ if $X$ is in $\und{Fib}_W S$, can we
\printedpage{121}
% head: A. Grothendieck : Dérivateurs, Ch. VI, Édition des Universités de Montpellier II et Paris VII
% foot: Version 9 novembre 2001
\msp{page 196}
hope to get by with the weakened variant of W(4b) obtained by restricting ourselves in the statement of W(4b) (p. 171) to the case where $X$, $Y$ are in $\und{Parf}_{W''_\omega} S$. This time, that does indeed suffice to imply that the functor
\[
\lambda : (\und{Fib}_W S)(W_S^{\mathrm{f}})^{-1} \longrightarrow (\und{Fib}_W S)(W_S)^{-1}
\]
is an equivalence of categories, since we can ``compute'' it just as well using $\und{Parf}_{W''_\omega} S$ instead of $\und{Fib}_W S$, now the hypothesis W(4b($\mathrm{W}''_\omega$)) implies that on $\und{Parf}_{W''_\omega} S$ the two localizers $W_S^{\mathrm{f}}$, $W_S$ are \textit{equal.}

We need a usable practical conclusion from all these perplexities and observations. Let us try

\textbf{Proposition.} \textit{Suppose that we are under one of the following conditions.}

\begin{enumerate}[label={}]
\item[$1^o$)] \textit{The axiom} W(4b) \textit{is satisfied in the case of categories $X$, $Y$ which are in} $\und{Fibsc}_{0,W} S$ (\footnote[95]{but this is \textit{equivalent} to W(4b)! Cf. XVI 4.9. loc. cit.}) \textit{(this condition on $W$ is in fact self-dual).}

\item[$2^o$)] \textit{The axiom} W(4b) \textit{is satisfied in the case of categories $X$, $Y$ which are in} $\und{Parf}_{W''_\omega} S$\textit{, where $W''_\omega$ is the smallest set $\Sigma \subseteq$ Fl}(Cat) \textit{satisfying the conditions} $\mathrm{W}(1, 2\text{ bis}, 3, 6\text{ bis})$. \textit{(We can also make do with the condition, stronger but more understandable, that} W(4b) \textit{is OK if $X$, $Y$ are in} $\und{Parf}_W S$ (\footnote[96]{But this also implies W(4b), cf. loc. cit.}).\textit{) Moreover, we assume that $W$ satisfies} W(6 bis).
\end{enumerate}

\textit{Under these conditions, the corresponding [?] functor $\lambda$}
\[
\mathrm{HOT}^{\mathrm{lc}}_W(S) \longrightarrow \mathrm{HOT}^{\mathrm{lc!}}_W(S)
\]
\textit{is an equivalence of categories (cf. diagram p. 178), so $(\rho, i)$, $(\sigma, j)$ can be regarded as pairs of functors between two categories (cf. $(*)$ p. 181).} These are pairs of adjoint functors, \textit{so $i$, $j$ are fully faithful, and we have the adjunction formula $(*)$, page 182.}

\msp{page 197}
To prove the adjunction of $(\rho, i)$ (or dually of $(\sigma, j)$), we again express the two categories at hand $\mathrm{HOT}_W(S)$, $\mathrm{HOT}^{\mathrm{lc}}_W(S)$ as $(\mathrm{Cat}/S)(W_S^{\mathrm{d}})^{-1}$, $(\mathrm{Cat}/S)(W_S)^{-1}$, the functor $\rho$ being the canonical localization going from the first to the second, and the functor $i$ being given by an inversion functor
\[
\gamma : (\mathrm{Cat}/S, W_S) \longrightarrow (\mathrm{Cat}/S, \overline{W_S^{\mathrm{d}}})
\]
which is $C$ in case $1^o$, $\Phi_\infty$ in case $2^o$. The adjunction morphisms
\[
\rho i \longrightarrow \mathrm{id}, \qquad i\rho \longleftarrow \mathrm{id}
\]
should then be described as in the case where we simply assume W(4b), and the compatibilities are seen in the same way: we are reduced to proving that two arrows
\printedpage{122}
% head: A. Grothendieck : Dérivateurs, Ch. VI, Édition des Universités de Montpellier II et Paris VII
% foot: Version 9 novembre 2001
\[
\begin{tikzcd}
\gamma(X) \arrow[r, shift left, "i_{\gamma(X)}"] \arrow[r, shift right, "\gamma(i_X)"'] & \gamma(\gamma(X))
\end{tikzcd}
\]
are equal in $(\mathrm{Cat}/S)(W_S^{\mathrm{d}})^{-1}$. The verification in case $1^o$), where $\gamma = C$, should not be changed. In case $2^o$) where $\gamma = \Phi_\infty$, I have not done it. I will do it only under the pressure of need.

\subsection*{\circnum{6} Contravariant law of $\mathrm{HOT}_W(S)$, $\mathrm{HOT}^{\mathrm{lc}}_W(S)$.}

There is for
\[
f : S' \longrightarrow S
\]
a homomorphism from the diagram
\[
\begin{tikzcd}[column sep=small]
& \mathrm{HOT}^{\mathrm{lc}}_W(S') \arrow[dl] \arrow[dr] & & \text{\texttt{[Rather}} & & \mathrm{HOT}^{\mathrm{lc}}_W(S) \arrow[dl] \arrow[dr] & \\
\mathrm{HOT}_W(S') & & \mathrm{HOT}^{o}_W(S') & & \mathrm{HOT}_W(S) & & \mathrm{HOT}^{o}_W(S)\ \text{\texttt{]}}
\end{tikzcd}
\]
\msp{page 198}
into the similar diagram for $S$ \texttt{[rather} $S'$\texttt{]}, by functors $f^*$ which are defined respectively on
\[
\begin{aligned}
\und{Liss}_W S &\quad \text{for}\quad \mathrm{HOT}_W(S) & \quad\longrightarrow\quad & \mathrm{HOT}_W(S')\\
\und{Prop}_W S &\quad \text{for}\quad \mathrm{HOT}^{o}_W(S) & \quad\longrightarrow\quad & \mathrm{HOT}^{o}_W(S')\\
\und{Fib}_W S &\quad \text{for}\quad \mathrm{HOT}^{\mathrm{lc}}_W(S) & \quad\longrightarrow\quad & \mathrm{HOT}^{\mathrm{lc}}_W(S'),
\end{aligned}
\]
through the functors % [sic] intermédiare -> intermédiaire, missing i
\[
f^* : X \longmapsto X \times_S S' : \left\{
\begin{aligned}
\und{Liss}_W S &\quad\longrightarrow\quad \und{Liss}_W S'\\
\und{Prop}_W S &\quad\longrightarrow\quad \und{Prop}_W S'\\
\und{Fib}_W S &\quad\longrightarrow\quad \und{Fib}_W S',
\end{aligned}
\right.
\]
which functor is compatible with the relevant localizers, which are $W_S^{\mathrm{f}}$ in all cases. On the other hand, in the absence of draconian hypotheses on $W$ or on $S$, we cannot say that $f^*$ is compatible with $W_S$, $W_{S'}$ on any of the categories of the diagram p. 175, allowing us to define $\mathrm{HOT}^{\mathrm{lc!}}_W(S)$, so as to give us a functor % [sic] pretinents -> pertinents, transposed letters
\[
\mathrm{HOT}^{\mathrm{lc!}}_W(S) \longrightarrow \mathrm{HOT}^{\mathrm{lc!}}_W(S').
\]
In the absence of axioms on $W$, the only case known to me is that where $f$ is \textit{a universal $W$-equivalence,} i.e. a $W$-fibration with $W$-aspheric fibers, in which case it is obvious that for \textit{every} arrow $u : X \longrightarrow Y$ in $\mathrm{Cat}/S$, $u \in W_S$ if and only if $(u' : X' \longrightarrow Y') \in W_{S'}$. We then find a homomorphism from the complete diagram $(*)$ p. 176 \texttt{[rather p. 178]} for $S$, into the corresponding diagram for $S'$.

If we assume W(4b), then $W_S = W_S^{\mathrm{f}}$ on $\und{Fib}_W S$, so on that category the functor $f^*$ is compatible (no merit!) with the localizers $W_S$, $W_{S'}$, but this is not really new
\printedpage{123}
% head: A. Grothendieck : Dérivateurs, Ch. VI, Édition des Universités de Montpellier II et Paris VII
% foot: Version 9 novembre 2001
\msp{page 199}
compared with what was said above, for $\und{Fib}_W S$, $\und{Fib}_W S'$, and for the $\mathrm{HOT}^{\mathrm{lc}}_W$.

But if we assume W(4a), and $f \in \und{Fib}_W$, then $f^*$ on all of $\mathrm{Cat}/S$ is compatible with the localizers $W_S$, $W_{S'}$, and we find a computation of $f^*$ on $\mathrm{HOT}^{\mathrm{lc!}}_W(S) \approx \mathrm{HOT}^{\mathrm{lc}}_W(S)$ by
\[
f^*(\mathrm{hot}^{\mathrm{lc}}_{S,W}(X)) = \mathrm{hot}^{\mathrm{lc}}_{S',W}(X \times_S S').
\]

So this is a case, exceptionally, where the functors $f^*$ for the $\mathrm{HOT}_W$ are compatible not only with the inclusion functors $i$, $j$ (of page 176 \texttt{[rather p. 178]}), but also with the localization functors $\rho$, $\sigma$. % [sic] exceptionellement -> exceptionnellement, French spelling has double n

\textbf{Omission:} If $p : S' \longrightarrow S$ is $W$-proper, then $p^* : \mathrm{Cat}/S \longrightarrow \mathrm{Cat}/S'$ is compatible with $W_S^{\mathrm{d}}$, $W_{S'}^{\mathrm{d}}$, and $p^* : \mathrm{HOT}_W(S) \longrightarrow \mathrm{HOT}_W(S')$ can be ``computed'' by the formula
\[
p^*(\mathrm{hot}_{S,W}(X)) = \mathrm{hot}_{S',W}(X \times_S S')
\]
for \textit{every} $X$ in $\mathrm{Cat}/S$.

\subsection*{\circnum{7} Product in $\mathrm{HOT}_W(S)$.}

The functor
\[
(X, Y) \longmapsto X \times_S Y : \und{Liss}_W S \times \und{Liss}_W S \longrightarrow \und{Liss}_W S
\]
is compatible with the localizers $W_S^{\mathrm{f}} \times W_S^{\mathrm{f}}$ and $W_S^{\mathrm{f}}$, so passes to the localized categories as
\[
\left\{
\begin{aligned}
\mathrm{HOT}_W(S) \times \mathrm{HOT}_W(S) &\quad\longrightarrow\quad \mathrm{HOT}_W(S)\\
(\xi, \eta) &\quad\longmapsto\quad \xi \times_S \eta
\end{aligned}
\right.
\]
(\footnote[97]{\textbf{NB} If $X$, $Y$ are \textit{arbitrary} in $\mathrm{Cat}/S$, we have $\mathrm{hot}_{S,W}(X) \times_S \mathrm{hot}_{S,W}(Y) = \mathrm{hot}_{S,W}(\Phi_S(X) \times_S \Phi_S(Y))$. If $X$ or $Y$ is smooth over $S$, we can replace $\Phi(X)$ resp. $\Phi(Y)$ there by $X$ resp. $Y$.}). This product is associative, unital, commutes with all the usual compatibilities, as follows from the fact that that's how it is in $\und{Liss}_S$. The unit element is the class of $S \xrightarrow{\mathrm{id}} S$. The subcategory $\mathrm{HOT}^{\mathrm{lc}}_W(S)$ is stable under the product, which induces
\msp{page 200}
\[
\mathrm{HOT}^{\mathrm{lc}}_W(S) \times \mathrm{HOT}^{\mathrm{lc}}_W(S) \longrightarrow \mathrm{HOT}^{\mathrm{lc}}_W(S),
\]
itself induced by the product functor $(X, Y) \longrightarrow X \times_S Y$ in $\und{Fib}_W S$.

If $X$ is arbitrary in $\mathrm{Cat}/S$, the functor
\[
\mathrm{mult}_{X/S} : Y \longmapsto X \times_S Y : \und{Liss}_W S \longrightarrow \mathrm{Cat}/S
\]
is compatible with the localizers $W_S^{\mathrm{d}}$, for this functor can be written as the composite of the corresponding functor \textit{with values in $\und{Liss}_W X$,} which is compatible with the localizers % [sic] correspondent -> correspondant, adjective intended